\section{Block separation and intersection}\label{sec:block_separation_intersection}

\begin{lemma}[Block-diagonal commutant reduction from span extension]
    \label{lem:block_diagonal_commutant_reduction_from_projection_span}
    \notready
    \uses{def:block_projection, def:block_diagonal_matrix}
    Let $S$ be a family of virtual matrices. If each sector projection satisfies
    \[
        P_j\in\spn S
    \]
    and a boundary matrix $X$ satisfies
    \[
        XM=MX\qquad (M\in S),
    \]
    then, for every sector $j$,
    \[
        XP_j=P_jX,
    \]
    and hence
    \[
        P_jXP_k=0\qquad (j\ne k).
    \]
    Thus $X$ is block diagonal with respect to the sector decomposition.
\end{lemma}

\begin{proof}
    \notready
    \uses{lem:block_diagonal_of_commutes_block_projection,
        lem:block_diagonal_iff_off_block_zero}
    The equations $XM=MX$ are linear in $M$, so they hold for every
    $M\in\spn S$. Substituting $M=P_j$ gives $XP_j=P_jX$. Multiplying this
    identity on the left by $P_j$ and on the right by $P_k$ with $j\ne k$
    gives $P_jXP_k=0$.
\end{proof}

\begin{lemma}[Block projections from product-algebra span]
    \label{lem:block_projection_span_from_product_algebra_span}
    \notready
    Let $T_a = (T_a(i))_i$ be a family of tuples spanning the full product algebra
    $\prod_i \MN{n_i}$, and let $c_i\neq0$ be nonzero scalars. Then, for every
    sector $k$, the sector projection $P_k$ lies in the span of the scaled
    dependent block-diagonal matrices
    \[
        \bigoplus_i c_i T_a(i).
    \]
\end{lemma}

\begin{proof}
    \notready
    Apply the linear map $(M_i)_i \mapsto \bigoplus_i c_i M_i$ to the full-span
    product-algebra element with $(c_k)^{-1} I$ in the $k$-th component and zero in
    every other component. Its image is exactly $P_k$.
\end{proof}

\begin{lemma}[Sector projections from blockwise word span]
    \label{lem:block_projection_span_from_product_word_span}
    \notready
    If the simultaneous length-$m$ block-word tuples
    \[
        \omega \longmapsto (A_j^{\omega})_j
    \]
    span the full product algebra $\prod_j \MN{D_j}$ and every $\mu_j$ is nonzero, then the
    pulled-back length-$m$ word products of $\bigoplus_j\mu_j A_j$ span every
    virtual sector projection $P_j$. Consequently, any boundary matrix commuting with all
    those direct-sum words has block-diagonal pulled-back form, and its off-block entries vanish.
\end{lemma}

\begin{proof}
    \notready
    \uses{lem:block_projection_span_from_product_algebra_span,
        lem:block_diagonal_commutant_reduction_from_projection_span}
    First embed the product algebra linearly as dependent block-diagonal matrices, scaling the
    $j$-th block by $\mu_j^m$. Since $\mu_j^m\neq0$, the product-algebra span contains the
    tuple with $(\mu_j^m)^{-1}I$ in one component and zero elsewhere; its scaled diagonal image is
    exactly $P_j$. The word-evaluation formula for the direct-sum tensor identifies these scaled
    diagonal matrices with the pulled-back direct-sum word products, and the commutant conclusion
    follows from Lemma~\ref{lem:block_diagonal_commutant_reduction_from_projection_span}.
\end{proof}

\begin{definition}[Pairwise block-separating equations]
    \label{def:pairwise_block_separating_words}
    \lean{MPSTensor.HasBlockSelectorOn}
    \lean{MPSTensor.HasPairBlockSeparatingWords}
    \leanok
    Fix a finite target set $T$ of block indices. A length-$S$ block-separating
    equation for $k$ against $T$ is a choice of coefficients $c_\omega$ such that
    \[
        \sum_{|\omega|=S} c_\omega A_\ell^\omega
        =
        \begin{cases}
            I, & \ell=k,\\
            0, & \ell\in T.
        \end{cases}
    \]
    The pairwise version requires this equation for each ordered pair of
    distinct blocks, with $T=\{j\}$.
\end{definition}

\begin{definition}[Pair trace-separation criterion]
    \label{def:pair_trace_separation_words}
    \lean{MPSTensor.pairAlgSpan}
    \lean{MPSTensor.pairWordTuple}
    \lean{MPSTensor.PairWordTupleSpanTop}
    \lean{MPSTensor.PairTraceSeparatingAt}
    \lean{MPSTensor.pairEvalWordTuple}
    \lean{MPSTensor.pairCumulativeSpan}
    \lean{MPSTensor.PairCumulativeWordTupleSpanTop}
    \lean{MPSTensor.PairTraceSeparatingUpTo}
    \lean{MPSTensor.PairTraceSeparatingAll}
    \lean{MPSTensor.pairAllWordsSpan}
    \lean{MPSTensor.PairAllWordsSpanTop}
    \leanok
    For two blocks $A$ and $B$, the length-$S$ pair word tuple is
    $w \mapsto (A^w,B^w)$. The pair product-span condition says that these
    tuples span $\MN{D_A}\times\MN{D_B}$. Equivalently, in the trace-dual
    formulation, the only pair of test matrices $(\Delta_A,\Delta_B)$ whose
    trace pairing with every length-$S$ tuple vanishes is $(0,0)$. The
    cumulative variant replaces one fixed length by all word lengths at most
    $S$, and the all-length variant allows every finite word.
\end{definition}

\begin{lemma}[Three-block trace reduction from a nonzero middle matrix]
    \label{lem:three_block_trace_reduction_nonzero_middle}
    \lean{MPSTensor.exists_right_trace_test_of_three_block_trace_relation_left}
    \lean{MPSTensor.exists_left_trace_test_of_three_block_trace_relation_right}
    \leanok
    \uses{def:l_blk_injective, lem:block_injective_two_sided_word_span}
    Let $A$ be $L$-block injective, and let $\Delta_A\neq 0$.
    Suppose that, for all words $u,v,t$ of length $L$,
    \[
        \tr\!(\Delta_A A_v A_t A_u)
        + \tr\!(\Delta_B B_v B_t B_u)=0.
    \]
    Then for every $Z\in M_{D_A}(\mathbb C)$ there is a matrix
    $W\in M_{D_B}(\mathbb C)$ such that, for all words $t$ of
    length $L$,
    \[
        \tr(ZA_t)+\tr(WB_t)=0.
    \]
    This is the trace-dual algebraic step in the two-block case of the
    direct-sum lemma of \cite{PerezGarcia2007Matrix}.
\end{lemma}

\begin{proof}\leanok\uses{lem:block_injective_two_sided_word_span}
    Since $\Delta_A\neq0$, Lemma~\ref{lem:block_injective_two_sided_word_span}
    writes $Z$ as a linear combination of matrices $A_u\Delta_A A_v$.
    For one such term, cyclicity of the trace rewrites the three-block relation
    as
    \[
        \tr(A_u\Delta_A A_vA_t)
        +\tr(B_u\Delta_B B_vB_t)=0.
    \]
    Taking the same linear combination of the matrices $B_u\Delta_BB_v$
    gives the required $W$, and linearity gives the general case.
\end{proof}

\begin{lemma}[Finite-chain image inclusion from a three-block relation]
    \label{lem:three_block_relation_ground_space_inclusion}
    \lean{MPSTensor.leftTraceWordMap}
    \lean{MPSTensor.groundSpaceMap_eq_leftTraceWordMap}
    \lean{MPSTensor.groundSpace_eq_leftTraceWordMap_range}
    \lean{MPSTensor.leftTraceWordMap_range_le_of_three_block_trace_relation_left}
    \lean{MPSTensor.leftTraceWordMap_range_eq_of_three_block_trace_relation}
    \lean{MPSTensor.groundSpace_le_of_three_block_trace_relation_left}
    \lean{MPSTensor.groundSpace_eq_of_three_block_trace_relation}
    \leanok
    \uses{lem:three_block_trace_reduction_nonzero_middle}
    Under the hypotheses of
    Lemma~\ref{lem:three_block_trace_reduction_nonzero_middle}, the length-$L$
    spaces \(\mathcal G_L^A\) and \(\mathcal G_L^B\) satisfy
    \[
        \mathcal G_L^A \subseteq \mathcal G_L^B.
    \]
    If the corresponding nonzero middle matrix exists on both sides, then
    $\mathcal G_L^A=\mathcal G_L^B$.
\end{lemma}

\begin{proof}\leanok
    \uses{lem:three_block_trace_reduction_nonzero_middle}
    Write a vector in $\mathcal G_L^A$ as
    $t\mapsto \tr(A_tZ)$. The preceding lemma supplies $W$ with
    $\tr(ZA_t)+\tr(WB_t)=0$ for every word $t$. Cyclicity of the
    trace gives
    \[
        \tr(A_tZ)=\tr(ZA_t)=-\tr(WB_t)=\tr(B_t(-W)),
    \]
    so the vector lies in $\mathcal G_L^B$. The equality statement follows
    by applying the same argument with the two blocks exchanged.
\end{proof}

\begin{lemma}[Dimension step for two-block local spaces]
    \label{lem:direct_sum_two_block_dimension_step}
    \lean{MPSTensor.leftTraceWordMap_injective_of_isNBlkInjective}
    \lean{MPSTensor.leftTraceWordMap_range_finrank_eq_of_isNBlkInjective}
    \lean{MPSTensor.leftTraceWordMap_range_eq_of_range_le_of_isNBlkInjective_of_dim_ge}
    \lean{MPSTensor.leftTraceWordMap_range_eq_of_three_block_trace_relation_left_of_dim_ge}
    \lean{MPSTensor.not_three_block_trace_relation_left_of_isNBlkInjective_of_dim_gt}
    \lean{MPSTensor.groundSpace_finrank_eq_of_isNBlkInjective}
    \lean{MPSTensor.bondDim_eq_of_groundSpace_eq_of_isNBlkInjective}
    \lean{MPSTensor.bondDim_eq_and_groundSpace_eq_of_groundSpace_le_of_isNBlkInjective_of_dim_ge}
    \lean{MPSTensor.bondDim_eq_and_groundSpace_eq_of_three_block_trace_relation_left_of_dim_ge}
    \lean{MPSTensor.groundSpace_inf_eq_bot_of_not_bondDim_eq_and_groundSpace_eq_of_dim_ge}
    \lean{MPSTensor.pairTraceSeparatingAt_threeBlock_of_isNBlkInjective_of_dim_gt}
    \lean{MPSTensor.pairTraceSeparatingAt_threeBlock_of_isNBlkInjective_of_dim_ne}
    \lean{MPSTensor.chainGroundSpace_eq_of_groundSpace_eq}
    \lean{MPSTensor.mpvSubmodule_ne_of_not_exists_mpv_eq_smul}
    \lean{MPSTensor.not_exists_mpv_eq_smul_of_not_exists_forall_mpv_eq_mul}
    \lean{MPSTensor.not_bondDim_eq_and_groundSpace_eq_of_mpvSubmodule_ne}
    \lean{MPSTensor.groundSpace_inf_eq_bot_of_exists_not_forall_mpv_eq_mul_of_dim_ge}
    \lean{MPSTensor.pairTraceSeparatingAt_of_groundSpace_inf_eq_bot_of_injective_groundSpaceMap}
    \lean{MPSTensor.pairTraceSeparatingAt_of_groundSpace_inf_eq_bot_of_isNBlkInjective}
    \lean{MPSTensor.pairTraceSeparatingAt_threeBlock_of_exists_not_forall_mpv_eq_mul_of_dim_ge}
    \lean{MPSTensor.groundSpace_inf_eq_bot_of_blocksNotGaugePhaseEquiv_same_dim_of_dim_ge}
    \lean{MPSTensor.pairTraceSeparatingAt_threeBlock_of_blocksNotGaugePhaseEquiv_same_dim_of_dim_ge}
    \lean{MPSTensor.forall_pairTraceSeparatingAt_threeBlock_of_blocksNotGaugePhaseEquiv}
    \lean{MPSTensor.exists_forall_pairTraceSeparatingAt_threeBlock_of_blocksNotGaugePhaseEquiv}
    \leanok
    \uses{lem:three_block_relation_ground_space_inclusion,
        thm:chain_ground_space_eq_mpv_blk}
    If $A$ and $B$ are length-$L$ block-injective, the map
    $Z\mapsto(t\mapsto\tr(A_tZ))$ has image dimension $D_A^2$, and
    similarly for $B$. Therefore an inclusion
    $\mathcal G_L^A\subseteq \mathcal G_L^B$, together with
    $D_B\le D_A$, forces $D_A=D_B$ and
    $\mathcal G_L^A=\mathcal G_L^B$. In particular, the three-block
    trace relation from
    Lemma~\ref{lem:three_block_relation_ground_space_inclusion} gives this
    equality whenever the ordering \(D_B\le D_A\) is available. If the
    inclusion comes from the strictly larger block, $D_B<D_A$, then the same
    dimension step already contradicts the existence of the three-block trace
    relation. More generally, if an external source hypothesis rules out the
    simultaneous conclusion $D_A=D_B$ and
    $\mathcal G_L^A=\mathcal G_L^B$, then the three-block local spaces
    $\mathcal G_{3L}^A$ and $\mathcal G_{3L}^B$ have zero intersection.
    With injectivity at the three-block length, the strict-size branch itself
    already gives homogeneous pair trace separation at length $3L$.
    One such source hypothesis is non-proportionality of the periodic MPS vectors
    at a sufficiently long chain length, equivalently distinctness of the
    corresponding spaces $\mathbb C\,V^{(N)}(A)$: equality
    \(\mathcal G_L^A=\mathcal G_L^B\) gives
    equality of the periodic-chain constraint spaces, and injective uniqueness
    identifies those spaces with the two one-dimensional MPS spans. For same-dimensional
    separated BNT blocks, the non-equivalence condition supplies such
    non-proportional MPS vectors at arbitrarily large chain lengths; the
    direct-sum injectivity hypotheses remain explicit. Zero intersection of
    the two three-block local spaces also gives homogeneous pair trace
    separation at that length once the boundary-to-chain maps at the
    three-block length are injective. Combining the equal-dimension and
    unequal-dimension branches over all ordered block pairs gives one common
    three-block trace-separation length, and hence the positive-length
    blockwise word span required by the block-separation argument. In the
    source theorem this span is obtained from finite-length injectivity for
    each block. The proved specialization records the stronger length-one span
    condition separately; under that extra hypothesis the same trace-separation
    argument starts at the corresponding one-block length. For an ordered pair
    of blocks, the trace-separation conclusion has the form
    \[
        \left[
        \forall |\omega|=S,\;
        \tr(\Delta_A A^\omega)+\tr(\Delta_B B^\omega)=0
        \right]
        \Longrightarrow
        \Delta_A=\Delta_B=0.
    \]
    Equivalently, the homogeneous pair span is full:
    \[
        \spn\{(A^\omega,B^\omega):|\omega|=S\}
        =
        \MN{D_A}\times\MN{D_B}.
    \]
    Concatenating words propagates this equality to positive multiples of
    \(S\), and hence to a finite period window for every ordered pair. Combining
    that window with the conditional period-window homogenization theorem gives
    a single length \(S_0\) such that
    \[
        \spn\{(A_1^\omega,\ldots,A_g^\omega):|\omega|=S_0\}
        =
        \prod_{j=1}^g \MN{D_j}.
    \]
    The same conclusion applies to the corresponding normal canonical-form
    hypotheses with a basis of normal tensors once one-site injectivity is
    supplied explicitly, since those hypotheses carry the irreducibility and
    normalization assumptions used in the direct-sum comparison.
\end{lemma}

\begin{proof}\leanok
    \uses{lem:three_block_relation_ground_space_inclusion}
    Block injectivity makes the boundary-to-chain map injective, so
    $\dim\mathcal G_L^A=D_A^2$ and $\dim\mathcal G_L^B=D_B^2$.
    Inclusion gives $D_A^2\le D_B^2$, while $D_B\le D_A$ gives the
    reverse inequality. Hence $D_A=D_B$, and equality of dimensions inside
    the inclusion gives \(\mathcal G_L^A=\mathcal G_L^B\). In the strict-size
    branch this contradicts $D_B<D_A$. For the conditional directness
    statement, a vector in the intersection of the two three-block local spaces
    has two boundary representations. If the first boundary matrix is zero,
    the vector is zero. Otherwise the difference of the two coefficient
    formulas is exactly the homogeneous three-block trace relation, so the
    dimension step gives the forbidden collapse. In the equal-size branch,
    equal local spaces impose the same cyclic local constraints on a
    periodic chain. The injective uniqueness theorem then identifies both
    cyclic ground spaces with the spans of their MPV states, contradicting
    non-proportionality of those states.
\end{proof}

\begin{lemma}[All-length pair trace separation has finite cumulative length]
    \label{lem:all_length_pair_trace_separation_finite_cutoff}
    \lean{MPSTensor.pairEvalWordTuple_mem_pairCumulativeSpan}
    \lean{MPSTensor.pairCumulativeSpan_mono}
    \lean{MPSTensor.PairTraceSeparatingUpTo.mono}
    \lean{MPSTensor.pairTraceSeparatingUpTo_of_pairTraceSeparatingAt}
    \lean{MPSTensor.pairTraceSeparatingAt_of_pairWordTupleSpanTop}
    \lean{MPSTensor.pairWordTupleSpanTop_iff_pairTraceSeparatingAt}
    \lean{MPSTensor.pairCumulativeWordTupleSpanTop_of_pairTraceSeparatingUpTo}
    \lean{MPSTensor.pairTraceSeparatingUpTo_of_pairCumulativeWordTupleSpanTop}
    \lean{MPSTensor.pairCumulativeWordTupleSpanTop_iff_pairTraceSeparatingUpTo}
    \lean{MPSTensor.pairCumulativeWordTupleSpanTop_of_pairWordTupleSpanTop}
    \lean{MPSTensor.pairCumulativeWordTupleSpanTop_of_pairTraceSeparatingAt}
    \lean{MPSTensor.pairAllWordsSpanTop_of_pairTraceSeparatingAll}
    \lean{MPSTensor.pairTraceSeparatingAll_of_pairAllWordsSpanTop}
    \lean{MPSTensor.pairAllWordsSpanTop_iff_pairTraceSeparatingAll}
    \lean{MPSTensor.exists_pairCumulativeWordTupleSpanTop_of_pairAllWordsSpanTop}
    \lean{MPSTensor.exists_pairTraceSeparatingUpTo_of_pairTraceSeparatingAll}
    \lean{MPSTensor.exists_pairTraceSeparatingUpTo_iff_pairTraceSeparatingAll}
    \lean{MPSTensor.exists_forall_pairTraceSeparatingUpTo_of_forall_pairTraceSeparatingAll}
    \leanok
    \uses{def:pair_trace_separation_words}
    If no nonzero trace functional vanishes on all finite pair words, then
    there is a finite $S$ such that vanishing on every word of length at most
    $S$ already forces the test pair to be zero. For a finite family of ordered
    block pairs, these bounds can be replaced by one common bound.
\end{lemma}

\begin{proof}
    \leanok
    The trace condition is dual to the span of all finite pair words. If that
    span is the whole product algebra, every vector lies in some cumulative
    span. Since the product algebra is finite-dimensional, the increasing
    sequence of cumulative spans stabilizes; at the stabilization index it is
    already the whole product algebra. Applying the finite trace-duality
    criterion gives the stated finite bound. This proves the cumulative
    finite-dimensional step; the remaining BNT step is to convert it to a
    single homogeneous length. Cumulative span-top alone is not identity
    padding, because the length-zero word already contributes the identity.
\end{proof}

\begin{lemma}[Pair algebra projections are full]
    \label{lem:pair_algebra_projection_fullness}
    \lean{MPSTensor.pairAlgSpan_map_fst_eq_top_of_isInjective}
    \lean{MPSTensor.pairAlgSpan_map_snd_eq_top_of_isInjective}
    \lean{MPSTensor.pairEvalWordTuple_mem_pairAlgSpan}
    \lean{MPSTensor.pair_mul_mem_pairAllWordsSpan}
    \lean{MPSTensor.pairAlgSpan_toSubmodule_le_pairAllWordsSpan}
    \lean{MPSTensor.pairAllWordsSpanTop_of_pairAlgSpan_eq_top}
    \lean{MPSTensor.pairTraceSeparatingAll_of_pairAlgSpan_eq_top}
    \lean{MPSTensor.gaugePhaseEquiv_of_pairAlgSpan_eq_graph_algEquiv}
    \lean{MPSTensor.not_pairAlgSpan_eq_graph_algEquiv_of_not_gaugePhaseEquiv}
    \lean{MPSTensor.subdirect_matrix_pair_eq_top_of_dim_ne}
    \lean{MPSTensor.pairAlgSpan_eq_top_of_injective_dim_ne}
    \lean{MPSTensor.matrixPairSubalgebra_eq_top_of_axes}
    \lean{MPSTensor.matrixPairSubalgebra_eq_top_of_leftAxisIdeal_eq_top}
    \lean{MPSTensor.matrixPairSubalgebra_eq_top_of_rightAxisIdeal_eq_top}
    \lean{MPSTensor.matrixPairSubalgebra_eq_graph_algEquiv_of_axisIdeals_eq_bot}
    \lean{MPSTensor.subdirect_matrix_pair_eq_top_or_eq_graph_algEquiv}
    \lean{MPSTensor.pairTraceSeparatingAll_of_injective_dim_ne}
    \leanok
    \uses{def:pair_trace_separation_words}
    If one of the two blocks is injective, then the corresponding coordinate
    projection of the pair algebra generated by the simultaneous one-letter
    pairs is the full matrix algebra. Moreover, every finite pair word lies in
    this pair algebra, and the all-word linear span is the submodule generated
    by the pair algebra. If the pair algebra is the graph of an algebra
    automorphism of the matrix algebra, Skolem--Noether turns this graph case
    into gauge-phase equivalence, so under non-gauge-equivalence the graph
    branch is impossible. For two blocks of unequal bond dimensions, the graph
    branch is also impossible because it would identify matrix algebras with
    different finite dimensions over $\C$.
\end{lemma}

\begin{proof}
    \leanok
    Projecting the pair generators to one coordinate gives exactly the
    one-block generator family. Injectivity says that the linear span of this
    family is the full matrix algebra, and the linear span is contained in the
    unital algebra it generates. The pair-word membership follows by induction
    on the word, using closure of the generated algebra under multiplication.
    Conversely, the all-word linear span is closed under componentwise
    multiplication, so it contains the algebra generated by the pair generators.
    In the graph case, each generator satisfies $B_i=\varphi(A_i)$; applying
    Skolem--Noether to $\varphi$ gives an ordinary gauge, hence a gauge-phase
    equivalence with scalar $1$.
\end{proof}

\begin{lemma}[Conditional padding from exact homogeneous identity witnesses]
    \label{lem:pair_trace_separation_homogenization_with_padding}
    \lean{MPSTensor.pairEvalWordTuple_mem_span_pairWordTuple_length}
    \lean{MPSTensor.pair_mul_mem_span_pairWordTuple_add}
    \lean{MPSTensor.pairIdentity_mem_pairWordTupleSpan_zero}
    \lean{MPSTensor.pairIdentity_mem_pairWordTupleSpan_mul}
    \lean{MPSTensor.pairIdentity_mem_pairWordTupleSpan_add}
    \lean{MPSTensor.pairIdentity_mem_pairWordTupleSpan_add_mul}
    \lean{MPSTensor.pairIdentity_mem_pairWordTupleSpan_eventually_of_period_window}
    \lean{MPSTensor.pairIdentity_mem_pairWordTupleSpan_eventually_of_pairWordTupleSpanTop_period_window}
    \lean{MPSTensor.pairIdentity_mem_pairWordTupleSpan_eventually_of_pairAllWordsSpanTop_period_window}
    \lean{MPSTensor.pairWordTupleSpanTop_of_pairCumulativeWordTupleSpanTop_of_identity_padding}
    \lean{MPSTensor.pairTraceSeparatingAt_of_pairTraceSeparatingUpTo_of_identity_padding}
    \lean{MPSTensor.exists_pairTraceSeparatingAt_of_pairTraceSeparatingUpTo_of_eventual_identity_padding}
    \lean{MPSTensor.exists_pairTraceSeparatingAt_of_pairAllWordsSpanTop_of_identity_padding}
    \lean{MPSTensor.exists_forall_pairTraceSeparatingAt_of_pairTraceSeparatingAll_of_identity_padding}
    \lean{MPSTensor.exists_forall_pairTraceSeparatingAt_of_pairTraceSeparatingAll_of_identity_period_windows}
    \lean{MPSTensor.exists_forall_pairTraceSeparatingAt_of_pairSpanTop_period_windows}
    \lean{MPSTensor.exists_pairTraceSeparatingAt_of_not_gaugePhaseEquiv_of_pairWordTupleSpanTop_period_window}
    \lean{MPSTensor.exists_pairTraceSeparatingAt_of_injective_dim_ne_of_pairWordTupleSpanTop_period_window}
    \lean{MPSTensor.pairAlgSpan_eq_top_of_injective_not_gaugePhaseEquiv}
    \lean{MPSTensor.pairTraceSeparatingAll_of_injective_not_gaugePhaseEquiv}
    \lean{MPSTensor.pairTraceSeparatingAll_of_injective_not_gaugePhaseEquiv_cast_left}
    \lean{MPSTensor.pairTraceSeparatingAll_and_exists_pairTraceSeparatingUpTo_of_injective_not_gaugePhaseEquiv}
    \lean{MPSTensor.exists_pairTraceSeparatingAt_of_pairTraceSeparatingUpTo_of_identity_period_window}
    \leanok
    \uses{def:pair_trace_separation_words, lem:all_length_pair_trace_separation_finite_cutoff}
    Homogeneous identity-padding witnesses multiply by concatenation of words.
    Therefore, if one has a positive period and a full residue window of
    homogeneous pair spans, then the simultaneous identity pair lies in all
    sufficiently large homogeneous pair spans. Combined with all-length pair
    trace separation, this period-window input gives one homogeneous
    pair trace-separation length. This does not prove the block-separation
    span condition of
    \cite{PerezGarcia2007Matrix} by itself: fullness of spans up to a
    bounded length is weaker than the exact homogeneous input used here.
\end{lemma}

\begin{proof}
    \leanok
    The length-zero word gives the identity only at length zero, so it is not a
    positive-length padding theorem. The useful padding facts are homogeneous:
    if the identity pair is in the span at lengths $m$ and $n$, componentwise
    multiplication puts it in the span at length $m+n$. A period witness and a
    complete residue window therefore give homogeneous padding at every
    sufficiently large length, which is the input needed by the
    all-length-to-fixed-length implication.
\end{proof}

\begin{lemma}[Pair trace separation gives pairwise block separation]
    \label{lem:pair_trace_separation_to_pairwise_block_separation}
    \lean{MPSTensor.pairWordTupleSpanTop_of_pairTraceSeparatingAt}
    \lean{MPSTensor.hasBlockSelectorOn_of_pairWordTupleSpanTop}
    \lean{MPSTensor.hasBlockSelectorOn_of_pairTraceSeparatingAt}
    \lean{MPSTensor.hasPairBlockSeparatingWords_of_forall_pairWordTupleSpanTop}
    \lean{MPSTensor.hasPairBlockSeparatingWords_of_forall_pairTraceSeparatingAt}
    \lean{MPSTensor.exists_hasPairBlockSeparatingWords_of_exists_forall_pairTraceSeparatingAt}
    \leanok
    \uses{def:pairwise_block_separating_words, def:pair_trace_separation_words}
    At a fixed length $S$, the pair trace-separation criterion is dual to full
    pair product-span. Therefore a pair trace-separation witness supplies a
    linear combination of word evaluations that is identity on the first block
    and zero on the second.
    A common witness for all ordered distinct pairs gives pairwise
    block-separating words at the same length.
\end{lemma}

\begin{proof}
    \leanok
    If the pair span were proper, a nonzero linear functional would vanish on
    it. Nondegeneracy of the matrix trace pairing represents that functional
    as $(\Delta_A,\Delta_B)$, contradicting pair trace separation. Once the
    pair span is the whole product algebra, the tuple $(I,0)$ is in the span.
    Using the same coefficients on the ambient block family gives the required
    pairwise block-separating equation.
\end{proof}

\begin{lemma}[Block-separating equations from pairwise separation]
    \label{lem:block_separating_equations_from_pairwise_separation}
    \lean{MPSTensor.pointwise_mul_mem_span_wordTuple_add}
    \lean{MPSTensor.hasBlockSelectorOn_empty}
    \lean{MPSTensor.HasBlockSelectorOn.mul}
    \lean{MPSTensor.hasBlockSelectorOn_finset_of_pairBlockSeparatingWords}
    \lean{MPSTensor.hasBlockSelectorWords_of_forall_hasBlockSelectorOn_univ_erase}
    \lean{MPSTensor.hasBlockSelectorWords_of_pairBlockSeparatingWords}
    \lean{MPSTensor.propBlockInjective_of_common_blockInjective_of_pairBlockSeparatingWords}
    \lean{MPSTensor.wordTupleSpanTop_of_common_blockInjective_of_pairBlockSeparatingWords}
    \lean{MPSTensor.hasBiCF_of_common_blockInjective_of_pairBlockSeparatingWords}
    \leanok
    \uses{def:pairwise_block_separating_words}
    If every ordered pair of distinct blocks admits a common-length pairwise
    block-separating equation, then for each $k$ there are coefficients
    $c_\omega^{(k)}$ with
    \[
        \sum_{|\omega|=(g-1)S} c_\omega^{(k)} A_j^\omega
        =
        \begin{cases}
            I, & j=k,\\
            0, & j\ne k.
        \end{cases}
    \]
    Hence common block injectivity plus pairwise separation implies the
    finite-length blockwise physical--virtual correspondence in the
    block-injectivity proposition in \cite[lines~340--345]{Cirac2016MPDO_arXiv}
    and the block-injective canonical-form trace-separation condition.
\end{lemma}

\begin{proof}
    \leanok
    The empty word gives the identity tuple, corresponding to an empty target set.
    The span of word tuples is closed under pointwise multiplication because the
    product of a length-$L$ word and a length-$S$ word is the concatenated
    length-$(L+S)$ word. Inducting over the finite set of blocks different from
    $k$ and multiplying the corresponding pairwise equations keeps the value $I$
    on $k$ and introduces one additional zero block at each step. Extracting
    coefficients from the final tuple-span membership gives the coefficient form
    of the displayed block-separating equation.
\end{proof}

\begin{lemma}[Homogeneous padding for blockwise product spans]
    \label{lem:homogeneous_word_tuple_span_padding}
    \lean{MPSTensor.wordTupleSpanTop_add_of_identity_mem_span_wordTuple}
    \lean{MPSTensor.wordTupleSpanTop_add_of_wordTupleSpanTop}
    \leanok
    \uses{def:blockwise_product_word_span}
    Suppose the length-\(L\) simultaneous block-word tuples span
    \(\prod_j M_{D_j}(\mathbb C)\), and suppose the simultaneous identity tuple
    \((I,\ldots,I)\) lies in the span of the length-\(S\) simultaneous
    block-word tuples. Then the length-\(L+S\) simultaneous block-word tuples
    span \(\prod_j M_{D_j}(\mathbb C)\). In particular, two full homogeneous
    spans at lengths \(L\) and \(S\) give a full homogeneous span at length
    \(L+S\).
\end{lemma}

\begin{proof}
    \leanok
    Let \(M=(M_j)_j\) be a target tuple.  Write
    \[
        M_j=\sum_{|u|=L} a_u A^j_u,
        \qquad
        I=\sum_{|v|=S} b_v A^j_v
    \]
    simultaneously for every block \(j\).  Multiplying the two expansions gives
    \[
        M_j
        =
        M_j I
        =
        \sum_{|u|=L}\sum_{|v|=S} a_u b_v A^j_uA^j_v
        =
        \sum_{|u|=L}\sum_{|v|=S} a_u b_v A^j_{uv}.
    \]
    The words \(uv\) have length \(L+S\), so the target tuple lies in the
    length-\(L+S\) span.
\end{proof}

\begin{lemma}[Period-window propagation for blockwise product spans]
    \label{lem:period_window_word_tuple_span_padding}
    \lean{MPSTensor.wordTupleSpanTop_add_mul_of_wordTupleSpanTop}
    \lean{MPSTensor.wordTupleSpanTop_eventually_of_wordTupleSpanTop_period_window}
    \leanok
    \uses{def:blockwise_product_word_span,
        lem:homogeneous_word_tuple_span_padding}
    Let \(p>0\). Suppose the length-\(p\) simultaneous block-word tuples span
    \(\prod_j M_{D_j}(\mathbb C)\), and suppose that for every
    \(0\le r<p\), the length-\(s+r\) simultaneous block-word tuples span
    \(\prod_j M_{D_j}(\mathbb C)\). Then there is a length \(N\) such that,
    for every \(n\ge N\), the length-\(n\) simultaneous block-word tuples span
    \(\prod_j M_{D_j}(\mathbb C)\). More concretely, one may take \(N=s\):
    writing
    \[
        n=s+r+qp,\qquad 0\le r<p,
    \]
    the span at length \(s+r\) propagates to the span at length \(n\) by
    multiplying \(q\) times by the period-\(p\) span.
\end{lemma}

\begin{proof}
    \leanok
    The Euclidean division of \(n-s\) by \(p\) gives
    \(n=s+r+qp\) with \(0\le r<p\).  The assumed residue-window span gives the
    full product algebra at length \(s+r\). Applying homogeneous padding with
    the length-\(p\) span once gives the full product algebra at length
    \(s+r+p\); iterating gives the full product algebra at length \(s+r+qp\).
\end{proof}

\begin{lemma}[Blockwise word span from block-separating equations]
    \label{lem:product_word_span_from_bnt_block_separation}
    \lean{MPSTensor.wordTupleSpanTop_of_common_blockInjective_of_blockSelectorWords}
    \lean{MPSTensor.wordTupleSpanTop_of_common_blockInjective_of_pairBlockSeparatingWords}
    \leanok
    \uses{def:blockwise_product_word_span,
        lem:block_separating_equations_from_pairwise_separation,
        lem:homogeneous_word_tuple_span_padding}
    Let $J$ be a finite set of $r$ blocks, and let $(A_j)_{j\in J}$ be a
    finite family of tensors with common block injectivity at length $L$.
    Suppose there are length-$S$ block-separating equations: for each block
    $k$,
    \[
        \sum_{|\omega|=S} c_\omega^{(k)} A_j^\omega
        =
        \begin{cases}
            I, & j=k,\\
            0, & j\ne k.
        \end{cases}
    \]
    Then the simultaneous length-$(L+S)$ block-word evaluations span
    $\prod_j \MN{D_j}$. If the block-separating equations are first obtained
    from pairwise block-separating equations at length $S$, the same conclusion
    holds at length $L+(r-1)S$.
\end{lemma}

\begin{proof}
    \leanok
    \uses{lem:block_separating_equations_from_pairwise_separation}
    Block injectivity writes each target block matrix as a length-$L$ linear
    combination of word evaluations. For a target tuple $(M_j)_j$, write $M_k$
    in this form in block $k$, multiply by the block-separating equation for
    $k$, and sum over $k$. Thus, if
    $M_k=\sum_{|u|=L}a_u^{(k)}A_k^u$ and
    $\sum_{|v|=S}b_v^{(k)}A_j^v=\delta_{j,k}I$, then for every block $j$,
    \[
        M_j=
        \sum_{k\in J}\sum_{|u|=L}\sum_{|v|=S}
        a_u^{(k)}b_v^{(k)} A_j^u A_j^v .
    \]
    The right side is a linear combination of length-$(L+S)$ word tuples, so
    these tuples span the full product algebra.
    If one starts with pairwise separation instead, first use
    Lemma~\ref{lem:block_separating_equations_from_pairwise_separation} to
    multiply the pairwise equations into the full equations.
\end{proof}

\begin{lemma}[BNT block separation gives a finite blockwise product span]
    \label{lem:bnt_direct_sum_block_separation_word_span}
    \lean{MPSTensor.hasPairBlockSeparatingWords_threeBlock_of_blocksNotGaugePhaseEquiv_c1}
    \lean{MPSTensor.propBlockInjective_of_blocksNotGaugePhaseEquiv_directSum_c1}
    \lean{MPSTensor.wordTupleSpanTop_of_blocksNotGaugePhaseEquiv_directSum_selectors_c1}
    \lean{MPSTensor.exists_wordTupleSpanTop_of_blocksNotGaugePhaseEquiv_directSum_selectors_c1}
    \leanok
    \uses{lem:direct_sum_two_block_dimension_step,
        lem:pair_trace_separation_to_pairwise_block_separation,
        lem:block_separating_equations_from_pairwise_separation,
        lem:product_word_span_from_bnt_block_separation}
    Let \(A^1,\ldots,A^r\) be separated BNT blocks satisfying the
    irreducibility, left-canonicality, and normalized self-overlap hypotheses
    used in the direct-sum comparison. Fix \(L_0>0\). Assume that, for every
    block \(j\), the map
    \[
        \Gamma_{L_0}^{A^j}:M_{D_j}(\mathbb C)\to
        (\mathbb C^d)^{\otimes L_0},
        \qquad
        X\longmapsto
        \sum_{|u|=L_0}\tr(XA^j_u)\ket{u}
    \]
    is injective, and that the homogeneous word spans at the two source
    lengths are full:
    \[
        \operatorname{span}\{A^j_u:|u|=L_0+1\}=M_{D_j}(\mathbb C),
        \qquad
        \operatorname{span}\{A^j_v:|v|=3(L_0+1)\}=M_{D_j}(\mathbb C).
    \]
    Then for every ordered pair \(k\ne j\) there are coefficients
    \(c^{k,j}_\omega\) such that
    \[
        \sum_{|\omega|=3(L_0+1)} c^{k,j}_\omega A^k_\omega=I,
        \qquad
        \sum_{|\omega|=3(L_0+1)} c^{k,j}_\omega A^j_\omega=0.
    \]
    Consequently the simultaneous block-word tuples at length
    \[
        (L_0+1)+(r-1)3(L_0+1)
    \]
    span the full product algebra
    \[
        \prod_{j=1}^r M_{D_j}(\mathbb C).
    \]
\end{lemma}

\begin{proof}
    \leanok
    Lemma~\ref{lem:direct_sum_two_block_dimension_step} gives homogeneous
    trace separation at length \(3(L_0+1)\) for every ordered pair of distinct
    BNT blocks. Lemma~\ref{lem:pair_trace_separation_to_pairwise_block_separation}
    converts these trace-separation equations into the displayed pairwise
    block-separating equations. Multiplying the equations for the \(r-1\)
    blocks different from \(k\) gives a block-separating equation for \(k\), and
    Lemma~\ref{lem:product_word_span_from_bnt_block_separation} then supplies the
    displayed product span.
\end{proof}

\begin{lemma}[Blockwise matrix equality from trace pairings]
    \label{lem:block_matrices_equal_from_trace_pairings}
    \lean{MPSTensor.block_matrices_eq_of_wordTupleSpanTop_trace}
    \leanok
    \uses{def:blockwise_product_word_span}
    Let \(A^1,\ldots,A^g\) be block tensors whose length-\(m\) simultaneous word
    tuples span the product algebra. If, for every word \(w\) of length \(m\),
    two blockwise matrix families \(X_j\) and \(Y_j\) satisfy
    \[
        \sum_j\operatorname{tr}(X_jA^j_w)
        =
        \sum_j\operatorname{tr}(Y_jA^j_w)
    \]
    then \(X_j=Y_j\) for every block \(j\).
\end{lemma}

\begin{proof}
    \leanok
    \uses{def:blockwise_product_word_span}
    For every word \(w\) of length \(m\),
    \[
        \sum_j\operatorname{tr}\!\left((X_j-Y_j)A^j_w\right)
        =
        \sum_j\operatorname{tr}(X_jA^j_w)
        -
        \sum_j\operatorname{tr}(Y_jA^j_w)
        =
        0.
    \]
    Nondegeneracy of the product trace pairing, together with the spanning
    hypothesis, gives \(X_j-Y_j=0\) for every block \(j\).
\end{proof}

\begin{lemma}[Boundary identities from block-sum membership]
    \label{lem:blockwise_boundary_identities_from_left_trace_membership}
    \lean{MPSTensor.pgvwc07_boundary_identities_of_leftBoundaryComponent_mem_iSup}
    \leanok
    \uses{lem:block_matrices_equal_from_trace_pairings,
        def:blockwise_product_word_span,
        def:ground_space_parent}
    Let \(A^1,\ldots,A^g\) be block tensors whose length-\(m\) simultaneous word
    tuples span the product algebra. Let \(C^j_a\in M_{D_j}(\mathbb C)\) be
    matrices indexed by blocks \(j\) and physical letters \(a\). Suppose that,
    for every pair of boundary letters \(a,b\) and every word \(w\) of length
    \(m\), a vector \(\psi\) has the left-boundary trace form
    \[
        \psi(a,w,b)=
        \sum_j\operatorname{tr}\!\left(A^j_b C^j_a A^j_w\right)
    \]
    and if
    \[
        \psi\in\bigvee_j G_{m+2}(A^j),
    \]
    then there are matrices \(E_j\) such that, for every block \(j\) and all
    \(a,b\),
    \[
        A^j_bC^j_a=A^j_bE_jA^j_a
    \]
\end{lemma}

\begin{proof}
    \leanok
    \uses{lem:block_matrices_equal_from_trace_pairings,
        def:ground_space_parent,
        def:ground_space_map}
    Write \(\psi=\sum_j\phi_j\) with \(\phi_j\in G_{m+2}(A^j)\). For each
    \(j\), choose \(E_j\) whose image under the length-\(m+2\) ground-space map
    for \(A^j\) is \(\phi_j\). Comparing the coefficient of the word \(awb\)
    in the two expressions for \(\psi\) gives
    \[
        \sum_j\operatorname{tr}\!\left(A^j_b C^j_a A^j_w\right)
        =
        \sum_j\operatorname{tr}\!\left(A^j_bE_jA^j_aA^j_w\right).
    \]
    Lemma~\ref{lem:block_matrices_equal_from_trace_pairings} gives the displayed
    blockwise identities.
\end{proof}

\begin{lemma}[Blockwise boundary compatibility from traces]
    \label{lem:blockwise_boundary_compatibility_from_traces}
    \lean{MPSTensor.pgvwc07_blockwise_compatibility_of_trace_decomposition}
    \leanok
    \uses{lem:block_matrices_equal_from_trace_pairings}
    Let \(A^1,\ldots,A^g\) be block tensors with the following common
    word-span property: for a fixed \(m\), the simultaneous tuples
    \[
        w\longmapsto
        \left(A^1_w,\ldots,A^g_w\right),
        \qquad |w|=m,
    \]
    span the full product algebra
    \(\prod_j M_{D_j}(\mathbb C)\).
    Suppose that matrices \(C^j_a,D^j_b\in M_{D_j}(\mathbb C)\) satisfy,
    for all physical indices \(a,b\) and all words \(w\) of length \(m\),
    \[
        \sum_j \operatorname{tr}\!\left(A^j_b C^j_a A^j_w\right)
        =
        \sum_j \operatorname{tr}\!\left(D^j_b A^j_a A^j_w\right).
    \]
    Then
    \[
        A^j_bC^j_a=D^j_bA^j_a
    \]
    for every block \(j\) and all \(a,b\).
\end{lemma}

\begin{proof}
    \leanok
    \uses{lem:block_matrices_equal_from_trace_pairings}
    For fixed \(a,b\), set
    \[
        \Delta_j=A^j_bC^j_a-D^j_bA^j_a .
    \]
    The assumed trace equality says that
    \[
        \sum_j \operatorname{tr}\!\left(\Delta_jA^j_w\right)=0
    \]
    for every word \(w\) of length \(m\).  Since the tuples
    \((A^1_w,\ldots,A^g_w)\) span the product algebra, the product trace
    pairing is nondegenerate and hence every \(\Delta_j\) is zero.
\end{proof}

\begin{lemma}[Normalized boundary matrix identities]
    \label{lem:normalized_boundary_matrix_identities}
    \lean{MPSTensor.pgvwc07_boundary_matrix_identities_of_compatibility}
    \leanok
    Let $(A_a)_{a\in I}$, $(C_a)_{a\in I}$, and $(D_b)_{b\in I}$ be
    matrices in $M_D(\mathbb C)$, indexed by a finite physical alphabet. If
    \[
        \sum_a A_a A_a^\dagger=\Id,
        \qquad
        A_b C_a=D_b A_a
    \]
    for all $a,b\in I$, and if
    \[
        E=\sum_a C_a A_a^\dagger,
    \]
    then
    \[
        D_b=A_bE,
        \qquad
        A_bC_a=A_bEA_a
    \]
    for all $a,b\in I$.
\end{lemma}

\begin{proof}
    \leanok
    Multiply $D_b$ by the normalization:
    \[
        D_b
        =
        D_b\sum_a A_aA_a^\dagger
        =
        \sum_a D_bA_aA_a^\dagger
        =
        \sum_a A_bC_aA_a^\dagger
        =
        A_bE .
    \]
    Substituting this in $A_bC_a=D_bA_a$ gives
    $A_bC_a=A_bEA_a$.
\end{proof}

\begin{lemma}[Left-boundary summands are ground-space vectors]
    \label{lem:left_boundary_summands_mem_block_ground_spaces}
    \lean{MPSTensor.pgvwc07_sum_leftBoundaryComponents_mem_iSup_groundSpace}
    \leanok
    Let \(A^1,\ldots,A^g\) be block tensors, and let
    \(C^j_a,E^j\in M_{D_j}(\mathbb C)\) satisfy
    \[
        A^j_b C^j_a=A^j_bE^jA^j_a
    \]
    for every block \(j\) and all physical indices \(a,b\).
    For a word \(\sigma=(\sigma_1,\ldots,\sigma_{n+2})\), define
    \[
        \alpha_j(\sigma)
        =
        \operatorname{tr}\!\left(
            A^j_{\sigma_{n+2}} C^j_{\sigma_1}
            A^j_{\sigma_2}\cdots A^j_{\sigma_{n+1}}
        \right).
    \]
    Then
    \[
        \sum_j \alpha_j
        \in
        \bigvee_j G_{n+2}(A^j).
    \]
\end{lemma}

\begin{proof}
    \leanok
    Fix \(j\).  By the assumed boundary identity and cyclicity of the
    trace,
    \[
        \operatorname{tr}\!\left(
            A^j_{\sigma_{n+2}} C^j_{\sigma_1}
            A^j_{\sigma_2}\cdots A^j_{\sigma_{n+1}}
        \right)
        =
        \operatorname{tr}\!\left(
            A^j_{\sigma_1}A^j_{\sigma_2}\cdots
            A^j_{\sigma_{n+2}}E^j
        \right).
    \]
    Thus \(\alpha_j\) is the image of \(E^j\) under the parametrization
    of \(G_{n+2}(A^j)\).  Hence each \(\alpha_j\) lies in
    \(G_{n+2}(A^j)\), and the finite sum lies in the supremum of these
    subspaces.
\end{proof}

\begin{lemma}[Left-boundary trace decompositions lie in the block ground spaces]
    \label{lem:left_boundary_trace_decomposition_mem_block_ground_spaces}
    \lean{MPSTensor.pgvwc07_mem_iSup_groundSpace_of_trace_decomposition}
    \leanok
    Let \(A^1,\ldots,A^g\) be block tensors with common word-span separation at
    length \(n\), and suppose that
    \[
        \sum_a A^j_a A^{j\,\dagger}_a=\Id
    \]
    for every block \(j\).  Suppose a vector \(\psi\in(\mathbb C^d)^{\otimes(n+2)}\)
    has the left-boundary trace decomposition
    \[
        \psi=\sum_j\alpha_j,
        \qquad
        \alpha_j(i_1,\ldots,i_{n+2})
        =
        \operatorname{tr}\!\left(
            A^j_{i_{n+2}} C^j_{i_1}
            A^j_{i_2}\cdots A^j_{i_{n+1}}
        \right),
    \]
    and suppose that the two trace decompositions agree for every middle word:
    \[
        \sum_j
        \operatorname{tr}\!\left(
            A^j_b C^j_a A^j_w
        \right)
        =
        \sum_j
        \operatorname{tr}\!\left(
            D^j_b A^j_a A^j_w
        \right).
    \]
    Then
    \[
        \psi\in\bigvee_j G_{n+2}(A^j).
    \]
\end{lemma}

\begin{proof}
    \leanok
    The trace-decomposition equality and the common word-span hypothesis give
    \[
        A^j_bC^j_a=D^j_bA^j_a
    \]
    for every \(j,a,b\).  By the normalization,
    \[
        D^j_b
        =
        D^j_b\sum_a A^j_aA^{j\,\dagger}_a
        =
        \sum_a A^j_bC^j_aA^{j\,\dagger}_a
        =
        A^j_bE^j,
        \qquad
        E^j:=\sum_a C^j_aA^{j\,\dagger}_a.
    \]
    Hence \(A^j_bC^j_a=A^j_bE^jA^j_a\).  Lemma
    \ref{lem:left_boundary_summands_mem_block_ground_spaces} then gives
    \[
        \psi=\sum_j\alpha_j\in\bigvee_j G_{n+2}(A^j).
    \]
\end{proof}

\begin{lemma}[One-step block intersection from block restrictions]
    \label{lem:block_restrictions_mem_block_ground_spaces}
    \lean{MPSTensor.pgvwc07_mem_iSup_groundSpace_of_iSup_restrictions}
    \leanok
    \uses{lem:left_boundary_trace_decomposition_mem_block_ground_spaces}
    Let \(A^1,\ldots,A^g\) be block tensors with common word-span separation at
    length \(n\), and suppose
    \[
        \sum_a A^j_aA^{j\,\dagger}_a=\Id
    \]
    for every block \(j\).  If \(\psi\in(\mathbb C^d)^{\otimes(n+2)}\) satisfies
    \[
        \psi(a,-)\in\bigvee_jG_{n+1}(A^j),
        \qquad
        \psi(-,b)\in\bigvee_jG_{n+1}(A^j)
    \]
    for every physical index \(a,b\), then
    \[
        \psi\in\bigvee_jG_{n+2}(A^j).
    \]
\end{lemma}

\begin{proof}
    \leanok
    Choose finite decompositions of the fixed-boundary vectors:
    \[
        \psi(a,i_2,\ldots,i_{n+2})
        =
        \sum_j
        \operatorname{tr}\!\left(
            A^j_{i_2}\cdots A^j_{i_{n+2}}C^j_a
        \right),
    \]
    and
    \[
        \psi(i_1,\ldots,i_{n+1},b)
        =
        \sum_j
        \operatorname{tr}\!\left(
            A^j_{i_1}\cdots A^j_{i_{n+1}}D^j_b
        \right).
    \]
    Evaluating these two decompositions on
    \((a,w,b)\), where \(w\) has length \(n\), gives
    \[
        \sum_j\operatorname{tr}\!\left(A^j_bC^j_aA^j_w\right)
        =
        \sum_j\operatorname{tr}\!\left(D^j_bA^j_aA^j_w\right).
    \]
    The first decomposition also gives
    \[
        \psi=\sum_j\alpha_j,\qquad
        \alpha_j(i_1,\ldots,i_{n+2})
        =
        \operatorname{tr}\!\left(
            A^j_{i_{n+2}}C^j_{i_1}
            A^j_{i_2}\cdots A^j_{i_{n+1}}
        \right).
    \]
    Lemma~\ref{lem:left_boundary_trace_decomposition_mem_block_ground_spaces}
    therefore gives
    \[
        \psi\in\bigvee_jG_{n+2}(A^j).
    \]
\end{proof}

\begin{lemma}[One-step block intersection characterization]
    \label{lem:block_restrictions_iff_block_ground_spaces}
    \lean{MPSTensor.pgvwc07_mem_iSup_groundSpace_iff_iSup_restrictions}
    \leanok
    \uses{lem:block_restrictions_mem_block_ground_spaces}
    Under the hypotheses of
    Lemma~\ref{lem:block_restrictions_mem_block_ground_spaces}, an
    \((n+2)\)-site vector satisfies
    \[
        \psi\in\bigvee_jG_{n+2}(A^j)
    \]
    if and only if
    \[
        \psi(a,-)\in\bigvee_jG_{n+1}(A^j),
        \qquad
        \psi(-,b)\in\bigvee_jG_{n+1}(A^j)
    \]
    for every physical index \(a,b\).
\end{lemma}

\begin{proof}
    \leanok
    The reverse implication is
    Lemma~\ref{lem:block_restrictions_mem_block_ground_spaces}.  For the
    forward implication, write
    \[
        \psi=\sum_j\gamma_j,\qquad \gamma_j\in G_{n+2}(A^j).
    \]
    Each \(\gamma_j\) has fixed-boundary restrictions in \(G_{n+1}(A^j)\):
    \[
        \gamma_j(a,-)\in G_{n+1}(A^j),
        \qquad
        \gamma_j(-,b)\in G_{n+1}(A^j).
    \]
    Taking the finite sum over \(j\) gives the two restrictions of \(\psi\) in
    \(\bigvee_jG_{n+1}(A^j)\).
\end{proof}

\begin{lemma}[Subspace form of the one-step block intersection]
    \label{lem:block_restriction_intersection_subspace}
    \lean{MPSTensor.pgvwc07_iSup_groundSpace_eq_restriction_intersection}
    \leanok
    \uses{lem:block_restrictions_iff_block_ground_spaces}
    Under the hypotheses of
    Lemma~\ref{lem:block_restrictions_mem_block_ground_spaces}, set
    \[
        S_n:=\bigvee_jG_{n+1}(A^j).
    \]
    Then
    \[
        \left\{\psi:\forall b,\ \psi(-,b)\in S_n\right\}
        \cap
        \left\{\psi:\forall a,\ \psi(a,-)\in S_n\right\}
        =
        \bigvee_jG_{n+2}(A^j).
    \]
\end{lemma}

\begin{proof}
    \leanok
    Lemma~\ref{lem:block_restrictions_iff_block_ground_spaces} gives the
    displayed equivalence for each vector \(\psi\).  Interpreting the two sides
    of that equivalence as subspaces gives the stated equality.
\end{proof}

\begin{lemma}[Period-window form of the block intersection identity]
    \label{lem:period_window_block_restriction_intersection}
    \lean{MPSTensor.pgvwc07_iSup_restriction_intersection_eventually_of_period_window}
    \leanok
    \uses{lem:period_window_word_tuple_span_padding,
        lem:block_restriction_intersection_subspace}
    Let \(p>0\). Suppose the length-\(p\) simultaneous block-word tuples span
    \(\prod_j M_{D_j}(\mathbb C)\), and suppose that for every \(0\le r<p\),
    the length-\(s+r\) simultaneous block-word tuples span
    \(\prod_j M_{D_j}(\mathbb C)\). Assume also the normalization
    \[
        \sum_a A^j_aA^{j\,\dagger}_a=I
    \]
    for every block \(j\). Then there is a length \(N\) such that, for every
    \(n\ge N\), if
    \[
        S_n:=\bigvee_jG_{n+1}(A^j),
    \]
    then
    \[
        \left\{\psi:\forall b,\ \psi(-,b)\in S_n\right\}
        \cap
        \left\{\psi:\forall a,\ \psi(a,-)\in S_n\right\}
        =
        \bigvee_jG_{n+2}(A^j).
    \]
\end{lemma}

\begin{proof}
    \leanok
    Lemma~\ref{lem:period_window_word_tuple_span_padding} gives a length \(N\)
    such that the simultaneous block-word tuples span
    \(\prod_j M_{D_j}(\mathbb C)\) at every length \(n\ge N\). For each such
    \(n\), Lemma~\ref{lem:block_restriction_intersection_subspace} applies and
    gives the displayed subspace equality.
\end{proof}

\begin{lemma}[BNT block separation gives one block-intersection step]
    \label{lem:bnt_direct_sum_block_separation_block_intersection}
    \lean{MPSTensor.pgvwc07_iSup_restriction_intersection_of_bnt_directSum_selectors_c1}
    \leanok
    \uses{lem:bnt_direct_sum_block_separation_word_span,
        lem:block_restriction_intersection_subspace}
    Under the hypotheses of
    Lemma~\ref{lem:bnt_direct_sum_block_separation_word_span}, assume in addition the
    normalization
    \[
        \sum_a A^j_aA^{j\,\dagger}_a=I
    \]
    for every block \(j\). Let
    \[
        n=(L_0+1)+(r-1)3(L_0+1),
        \qquad
        S_n=\bigvee_jG_{n+1}(A^j).
    \]
    Then
    \[
        \left\{\psi:\forall b,\ \psi(-,b)\in S_n\right\}
        \cap
        \left\{\psi:\forall a,\ \psi(a,-)\in S_n\right\}
        =
        \bigvee_jG_{n+2}(A^j).
    \]
\end{lemma}

\begin{proof}
    \leanok
    Lemma~\ref{lem:bnt_direct_sum_block_separation_word_span} gives
    \[
        \operatorname{span}\{(A^1_w,\ldots,A^r_w): |w|=n\}
        =
        \prod_jM_{D_j}(\mathbb C).
    \]
    Substituting this span into
    Lemma~\ref{lem:block_restriction_intersection_subspace}, together with the
    normalization, gives the stated subspace equality.
\end{proof}

\begin{lemma}[BNT block separation with a block-injectivity window]
    \label{lem:bnt_direct_sum_block_separation_period_window_block_intersection}
    \lean{MPSTensor.pgvwc07_iSup_restriction_intersection_eventually_of_bnt_directSum_period_window}
    \leanok
    \uses{lem:direct_sum_two_block_dimension_step,
        lem:pair_trace_separation_to_pairwise_block_separation,
        lem:block_separating_equations_from_pairwise_separation,
        lem:period_window_block_restriction_intersection}
    Let \(A^1,\ldots,A^r\) be separated BNT blocks satisfying the
    irreducibility, left-canonicality, and normalized self-overlap hypotheses
    used in the direct-sum comparison. Assume \(L>1\), and assume that, for
    every block \(j\),
    \[
        \operatorname{span}\{A^j_a:a=0,\ldots,d-1\}=M_{D_j}(\mathbb C),
        \qquad
        \operatorname{span}\{A^j_u:|u|=L\}=M_{D_j}(\mathbb C),
        \qquad
        \operatorname{span}\{A^j_v:|v|=L+(L+L)\}=M_{D_j}(\mathbb C).
    \]
    Set
    \[
        S=L+(L+L),
        \qquad
        q=(r-1)S.
    \]
    Suppose that \(p>0\) and let \(s_0\) be a starting length. Assume that, for
    every block \(j\),
    \[
        \operatorname{span}\{A^j_u:|u|=p\}=M_{D_j}(\mathbb C),
    \]
    and that, for every \(0\le s<p+q\) and every block \(j\),
    \[
        \operatorname{span}\{A^j_v:|v|=s_0+s\}=M_{D_j}(\mathbb C).
    \]
    Assume also the normalization
    \[
        \sum_a A^j_aA^{j\,\dagger}_a=I
    \]
    for every block \(j\). Then there is \(N\) such that, for every \(n\ge N\),
    if
    \[
        S_n:=\bigvee_jG_{n+1}(A^j),
    \]
    then
    \[
        \left\{\psi:\forall b,\ \psi(-,b)\in S_n\right\}
        \cap
        \left\{\psi:\forall a,\ \psi(a,-)\in S_n\right\}
        =
        \bigvee_jG_{n+2}(A^j).
    \]
\end{lemma}

\begin{proof}
    \leanok
    The direct-sum dimension step gives, for each ordered pair of distinct
    blocks, length-\(S\) equations which are the identity on the first block and
    zero on the second. Multiplying these equations over the \(r-1\) other
    blocks gives a block-separating equation of total length \(q\). Combining
    this equation with block injectivity at length \(p\) gives the full product
    span at length \(p+q\). Combining it with the assumed block-injectivity
    window gives the full product span at lengths \(s_0+q+s\), for
    \(0\le s<p+q\). The period-window form of the block-intersection identity
    then gives the displayed equality for all sufficiently large \(n\).
\end{proof}

\begin{lemma}[Normalized block-separating equations give all large product spans]
    \label{lem:unital_block_separating_product_span}
    \lean{MPSTensor.wordTupleSpanTop_of_ge_of_common_blockInjective_of_unital_of_pairBlockSeparatingWords}
    \lean{MPSTensor.wordTupleSpanTop_of_ge_of_bnt_directSum_unital_c1}
    \leanok
    \uses{def:pairwise_block_separating_words,
        lem:bnt_direct_sum_block_separation_word_span,
        thm:wordspan_propagation_unital}
    Let \(A^1,\ldots,A^r\) be blocks satisfying the normalization
    \[
        \sum_a A^j_aA^{j\,\dagger}_a=I
    \]
    and assume that every block is injective at length \(L\):
    \[
        \operatorname{span}\{A^j_u:|u|=L\}=M_{D_j}(\mathbb C).
    \]
    If there are block-separating equations of length \(S\), then, for every
    \(n\ge L+(r-1)S\),
    \[
        \operatorname{span}\{(A^1_w,\ldots,A^r_w):|w|=n\}
        =
        \prod_j M_{D_j}(\mathbb C).
    \]
    In particular, for separated BNT blocks satisfying the same normalization,
    if \(L_0>0\) and \(\Gamma_{L_0}^{A^j}\) is injective for every block \(j\),
    the same conclusion holds with
    \[
        S=3(L_0+1),
        \qquad
        n\ge (L_0+1)+(r-1)3(L_0+1).
    \]
\end{lemma}

\begin{proof}
    \leanok
    By Theorem~\ref{thm:wordspan_propagation_unital}, the normalization
    propagates the length-\(L\) span equality to every larger homogeneous length:
    \[
        m\ge L
        \quad\Longrightarrow\quad
        \operatorname{span}\{A^j_u:|u|=m\}=M_{D_j}(\mathbb C).
    \]
    For \(n\ge L+(r-1)S\), take \(m=n-(r-1)S\). Combining the length-\(m\)
    block span with the \(r-1\) block-separating equations gives the displayed
    product span at length \(n\). In the BNT case, the injectivity of
    \(\Gamma_{L_0}^{A^j}\) gives the full span at length \(L_0\), and the
    normalization propagates it to the lengths \(L_0+1\) and \(3(L_0+1)\).
    Lemma~\ref{lem:bnt_direct_sum_block_separation_word_span} then supplies the
    block-separating equations with \(S=3(L_0+1)\).
\end{proof}

\begin{lemma}[Product span makes the sum of local spaces direct]
    \label{lem:product_span_block_image_direct_sum}
    \lean{MPSTensor.groundSpace_iSupIndep_of_wordTupleSpanTop}
    \lean{MPSTensor.groundSpace_iSupIndep_of_ge_of_bnt_directSum_unital_c1}
    \leanok
    \uses{def:ground_space_parent, lem:unital_block_separating_product_span}
    Suppose
    \[
        \operatorname{span}\{(A^1_w,\ldots,A^r_w):|w|=n\}
        =
        \prod_jM_{D_j}(\mathbb C).
    \]
    Then the sum of the local spaces \(G_n(A^j)\) is direct: if
    \[
        \phi_j\in G_n(A^j),
        \qquad
        \sum_j\phi_j=0,
    \]
    then \(\phi_j=0\) for every \(j\).  In particular, for the normalized BNT
    hypotheses above, this directness holds for every
    \[
        n\ge (L_0+1)+(r-1)3(L_0+1).
    \]
\end{lemma}

\begin{proof}
    \leanok
    Write
    \[
        \phi_j(\sigma)=\tr(A^j_\sigma X_j).
    \]
    Evaluating \(\sum_j\phi_j=0\) on each word \(w\) of length \(n\) gives
    \[
        \sum_j\tr(A^j_wX_j)=0.
    \]
    Since \(\tr(A^j_wX_j)=\tr(X_jA^j_w)\), the product-span equation gives, for
    every \((M_1,\ldots,M_r)\in\prod_jM_{D_j}(\mathbb C)\),
    \[
        \sum_j\tr(X_jM_j)=0.
    \]
    Taking only the \(j\)-th component nonzero and using nondegeneracy of the
    trace pairing yields \(X_j=0\), hence \(\phi_j=0\). The BNT statement follows
    by substituting Lemma~\ref{lem:unital_block_separating_product_span}.
\end{proof}
